% =============================================================
\section{Results}
\label{sec:results}
% =============================================================

\subsection{Persistence reference}

Table~\ref{tab:persistence} reports the persistence $\rmseday$ at each site and
horizon, which defines the denominator of $\skillday$ for all models.

\begin{table}[t]
  \centering
  \caption{Persistence test-set $\rmseday$ (W/m²). All deep-learning and
           statistical models are compared against these values.}
  \label{tab:persistence}
  \begin{tabular}{lccc}
    \toprule
    \textbf{Site} & \textbf{1h} & \textbf{3h} & \textbf{6h} \\
    \midrule
    El Paso   & 204.4 & 411.0 & 565.3 \\
    Uniandes  & 293.9 & 405.0 & 470.9 \\
    \bottomrule
  \end{tabular}
\end{table}

The 1-hour persistence error at El~Paso (204.4~W/m²) is notably lower than at
Uniandes (293.9~W/m²), reflecting El~Paso's more stable clear-sky regime where
short-horizon irradiance changes are small. At 6~hours, both sites have comparable
absolute difficulty (565.3 vs.\ 470.9~W/m²).

\subsection{Main results}

Table~\ref{tab:main} reports $\rmseday$ and $\skillday$ for all models on the
test set. Deep-learning results are mean~$\pm$~std over four seeds; statistical
models use a single run.

\begin{table*}[t]
  \centering
  \caption{Test-set results. $\rmseday$ in W/m². $\skillday$ relative to persistence
           (higher is better). DL models: mean~$\pm$~std over 4 seeds.
           Best $\skillday$ per site-horizon highlighted in \textbf{bold}.}
  \label{tab:main}
  \setlength{\tabcolsep}{5pt}
  \begin{tabular}{llcccccc}
    \toprule
    & & \multicolumn{3}{c}{\textbf{El Paso (César)}} &
         \multicolumn{3}{c}{\textbf{Uniandes (Bogotá)}} \\
    \cmidrule(lr){3-5}\cmidrule(lr){6-8}
    \textbf{Model} & \textbf{Metric} & \textbf{1h} & \textbf{3h} & \textbf{6h}
                                     & \textbf{1h} & \textbf{3h} & \textbf{6h} \\
    \midrule
    \multirow{2}{*}{Persistence}
      & $\rmseday$ & 204.4 & 411.0 & 565.3  & 293.9 & 405.0 & 470.9 \\
      & $\skillday$ & 0.000 & 0.000 & 0.000 & 0.000 & 0.000 & 0.000 \\
    \midrule
    \multirow{2}{*}{SARIMAX}
      & $\rmseday$ & 447.9 & 450.4 & 445.8 & 269.3 & 265.1 & 265.6 \\
      & $\skillday$ & $-1.205$ & $-0.093$ & $+0.210$ & $+0.081$ & $+0.353$ & $+0.435$ \\
    \midrule
    \multirow{2}{*}{ResNet-LSTM}
      & $\rmseday$ & $238.6 \pm 19.4$ & $238.5 \pm 36.9$ & $218.4 \pm 11.4$
                   & $246.7 \pm 3.6$ & $255.7 \pm 4.2$ & $297.9 \pm 7.6$ \\
      & $\skillday$ & $-0.167 \pm 0.095$ & $+0.420 \pm 0.090$ & $\mathbf{+0.614} \pm 0.020$
                    & $+0.161 \pm 0.012$ & $+0.369 \pm 0.010$ & $+0.368 \pm 0.016$ \\
    \midrule
    \multirow{2}{*}{GraphSAGE-LSTM}
      & $\rmseday$ & $244.7 \pm 7.0$ & $238.9 \pm 13.7$ & $\mathbf{206.1} \pm 6.4$
                   & $250.8 \pm 4.4$ & $252.5 \pm 1.3$ & $274.4 \pm 3.4$ \\
      & $\skillday$ & $-0.197 \pm 0.034$ & $+0.419 \pm 0.033$ & $\mathbf{+0.636} \pm 0.011$
                    & $+0.147 \pm 0.015$ & $\mathbf{+0.376} \pm 0.003$ & $+0.417 \pm 0.007$ \\
    \bottomrule
  \end{tabular}
\end{table*}

\subsection{Horizon-by-horizon analysis}

\paragraph{1-hour horizon.}
All deep-learning models fail to beat persistence at El~Paso: both ResNet-LSTM
($\skillday = -0.167 \pm 0.095$) and GraphSAGE-LSTM ($-0.197 \pm 0.034$) show
negative skill, meaning they produce higher RMSE than simply repeating the last
observation. SARIMAX is even worse ($-1.205$). At Uniandes the picture is more
nuanced: ResNet-LSTM achieves modest positive skill ($+0.161 \pm 0.012$) and
GraphSAGE-LSTM $+0.147 \pm 0.015$, while SARIMAX reaches $+0.081$. The key
difference between sites is the persistence reference: El~Paso's 1-hour persistence
RMSE (204.4~W/m²) is comparatively low because irradiance changes slowly under
predominantly clear-sky conditions, leaving little room for improvement from data-driven
methods. Uniandes' 1-hour persistence RMSE (293.9~W/m²) is higher due to the
frequent cloud passages, which create more structure for the models to exploit.

\paragraph{3-hour horizon.}
Both deep-learning models outperform persistence at El~Paso by approximately
42\% ($\skillday \approx 0.42$ for both architectures). The large standard deviation
for ResNet-LSTM ($\pm 0.090$) reflects seed-to-seed variability, with some seeds
achieving near-GraphSAGE performance and others considerably below. SARIMAX
exhibits slight negative skill ($-0.093$), confirming that linear time-series
models cannot exploit satellite spatial information.

At Uniandes, DL skill scores are $+0.369$ (ResNet) and $+0.376$ (GraphSAGE),
similar to each other and slightly below El~Paso despite the higher absolute
persistence error. SARIMAX shows strong skill at this horizon ($+0.353$),
approaching the DL models, which we attribute to the more regular temporal
autocorrelation structure of the Bogot{\'a} clear-sky index.

\paragraph{6-hour horizon.}
This is the most favourable horizon for the deep-learning models.
\textbf{GraphSAGE-LSTM achieves the best overall result}: $\rmseday = 206.1 \pm 6.4$~W/m²
and $\skillday = 0.636 \pm 0.011$ at El~Paso, reducing the persistence RMSE by
63.6\%. ResNet-LSTM is close ($0.614 \pm 0.020$) but consistently below GraphSAGE
across all seeds. SARIMAX achieves positive skill ($+0.210$) but is substantially
below the DL methods.

At Uniandes the DL advantage over SARIMAX narrows significantly: GraphSAGE-LSTM
achieves $0.417 \pm 0.007$ while SARIMAX reaches $0.435$, making SARIMAX the
best-performing model at this site-horizon combination. ResNet-LSTM underperforms
both ($0.368 \pm 0.016$), exhibiting the largest gap between the two sites.

\subsection{Statistical robustness}

The standard deviation across seeds provides a measure of training stability.
GraphSAGE-LSTM consistently shows lower seed-to-seed variability than ResNet-LSTM
(e.g., El~Paso 6h: $\pm 0.011$ vs.\ $\pm 0.020$), suggesting that the graph
structure regularises training and reduces sensitivity to weight initialisation.
The most unstable model is ResNet-LSTM at El~Paso 3h ($\pm 0.090$), which may
reflect the interaction between the high-irradiance regime and the CNN's sensitivity
to specific initialisation-driven features.
